\chapter{Sprint 3 : Couche d'intelligence}

\section{Introduction}
Le Sprint 3 est le sprint qui transforme Nova~AI en \emph{partenaire}
plutot qu'en \emph{outil de prise de rendez-vous}. A la fin du Sprint 2,
la plateforme repondait aux messages WhatsApp, tenait un calendrier,
faisait tourner un Kanban laboratoire, et alimentait une PWA patient.
Toutes ces fonctionnalites sont precieuses, mais un concurrent disposant
d'une equipe d'ingenierie similaire pourrait les repliquer en quelques
mois. Les fonctionnalites construites au Sprint 3 sont differentes :
elles transforment les donnees capturees en intelligence operationnelle
sur laquelle le proprietaire peut agir. C'est cette couche qui rend la
plateforme adherente.

Nous avons construit sept telles fonctionnalites au cours du sprint,
plus une (avatars vocaux) qui a ete implementee mais conservee derriere
un drapeau de fonctionnalite en attendant un abonnement ElevenLabs paye.
Le chapitre decrit chaque feature en suivant le meme modele : le
probleme metier qu'elle resout, le schema et le code qui la mettent en
oeuvre, le flux utilisateur, et les metriques operationnelles qu'elle
produit.

\section{Backlog du sprint}
\begin{itemize}
  \item Demand-fill auto-pinging.
  \item Apprentissage automatique des durees de consultation.
  \item Verificateur d'interactions medicamenteuses et d'allergies.
  \item Profils de vehicules pour garages.
  \item Surprises d'anniversaire IA.
  \item Check-in QR avec minuteur de consultation.
  \item File d'attente en temps reel pour laboratoires.
  \item Avatars vocaux (derriere un drapeau de fonctionnalite).
\end{itemize}

\section{Demand-Fill : remplissage automatique de creneaux}

\subsection{Le probleme metier}
Une clinique qui reserve, en moyenne, vingt rendez-vous par jour
remplit rarement la totalite. Les annulations laissent des trous. Un
mardi apres-midi avec cinq creneaux vides entre 14h00 et 17h00
represente un revenu reel perdu --- parfois l'equivalent d'une matinee
complete. Le proprietaire ne le remarque que retroactivement, quand la
journee est deja passee.

\subsection{La solution}
Une tache cron quotidienne (\texttt{demand-fill-runner} a 9h00) scanne
le calendrier de chaque entreprise pour les sept prochains jours.
Lorsqu'elle trouve un medecin avec trois creneaux contigus ou plus
entre 9h00 et 18h00 sur un jour ouvrable, elle cree un brouillon de
\emph{campagne}. La campagne inclut la date, le medecin, les heures
des creneaux ouverts, et une liste classee de patients qui
correspondent au service probable du creneau vide.

\subsection{Classement des candidats}
Le classement est implemente comme une fonction Postgres. La formule
de score favorise les patients qui n'ont pas visite depuis 30+ jours,
penalise ceux qui ont plusieurs no-shows, et recompense ceux qui ont
deja reserve le service candidat. Les patients pinges au cours des 14
derniers jours sont entierement exclus pour eviter le spam, ainsi que
ceux dont le drapeau de consentement marketing est false.

\subsection{Flux d'approbation par le proprietaire}
Le proprietaire voit les campagnes proposees dans un onglet
\emph{Demand Fill} de la barre laterale. Chaque carte de campagne
montre la date, le medecin, les heures de creneaux, et le compte de
candidats. Deux boutons : \emph{Voir les candidats} (ouvre une fenetre
modale avec la liste classee), \emph{Approuver et envoyer} (declenche
la diffusion WhatsApp).

\subsection{Auto-reservation WhatsApp}
Le client pinge recoit un message WhatsApp interactif :
\emph{\og Bonjour Mehdi, le Dr Bouazizi a des creneaux disponibles
le 12 mai : 14:30, 15:00, 15:30, 16:00. Voulez-vous en saisir un ? \fg{}}
avec deux boutons. Le webhook court-circuite le tap OUI vers une
reservation directe --- reclame le premier creneau restant, le retire
du tableau \texttt{slot\_times} de la campagne, incremente le compteur
\texttt{filled\_count}, et envoie une confirmation. Les conditions de
course sont gerees au niveau de la base : si deux patients tapent OUI
en meme temps, un seul obtient le creneau, l'autre recoit un message
poli \og ces creneaux viennent juste d'etre reserves \fg{}.

\subsection{KPI}
Le tableau de bord expose quatre tuiles KPI : campagnes envoyees, sept
derniers jours, creneaux remplis, taux de remplissage en pourcentage.
Apres un mois d'operation, un taux de remplissage typique se situe
entre 18 et 25\%.

\section{Apprentissage automatique des durees de consultation}

\subsection{Le probleme metier}
Une consultation programmee comme \og 30 minutes \fg{} prend rarement
30 minutes. Certains medecins courent systematiquement long, d'autres
court, et certains dependent du service : une consultation de suivi
avec un patient connu prend authentiquement dix minutes, tandis qu'une
consultation de premiere visite en cardiologie prend quarante-cinq
minutes. Quand le logiciel de reservation suppose 30 minutes pour
tout, la journee deborde.

\subsection{La solution}
Le schema a gagne des colonnes \texttt{arrived\_at} et
\texttt{completed\_at} sur \texttt{appointments}, plus une colonne
calculee generee \texttt{actual\_duration\_min}. Une vue
\texttt{v\_service\_actual\_durations} agrege les 90 derniers jours de
visites completees par (entreprise, medecin, service) : compte,
moyenne, mediane, p90.

Un cron hebdomadaire appelle une fonction Postgres
\texttt{duration\_suggestions} qui retourne la liste des services
deviant de plus de 15\% par rapport a leur duree programmee avec au
moins 5 echantillons. Pour chaque service, la suggestion est arrondie
au prochain creneau de 5 minutes.

\subsection{Opt-in d'application automatique}
Le proprietaire peut basculer un toggle \texttt{auto\_adjust\_duration}
par service. Quand le cron hebdomadaire s'execute, les services avec
le toggle actif voient leur \texttt{duration\_min} mis a jour
automatiquement ; les services sans le toggle recoivent une
notification avec la suggestion que le proprietaire peut appliquer
en un clic.

\subsection{Mecanique de capture}
Les durees reelles sont capturees de trois manieres : les boutons
\emph{Arrive} et \emph{Termine} du tableau de bord ; le flux de scan
QR decrit plus bas, qui horodate \texttt{arrived\_at} en effet
secondaire ; et un trigger base de donnees
\texttt{maybe\_mark\_arrived} qui ecoute les messages entrants
contenant des mots comme \og I'm here \fg{}, \og j'arrive \fg{},
\og wsalt \fg{} dans la demi-heure autour de l'horaire prevu.

\section{Verificateur d'interactions medicamenteuses et d'allergies}

\subsection{Le probleme metier}
Un medecin redigeant une prescription a besoin de savoir si le
medicament propose interagit avec ce que le patient prend deja, ou
s'il entre en conflit avec une allergie enregistree. Les systemes
hospitaliers de dossier medical electronique le font de routine ;
les cliniques privees en Tunisie ne le font presque jamais, parce
que le logiciel dedie de securite des medicaments est cher et
specifique a la locale. Une seule interaction manquee peut signifier
une hospitalisation. Sur une carriere, c'est le type d'erreur qui
termine un proces en faute professionnelle.

\subsection{La solution}
Le schema a gagne trois colonnes JSONB sur \texttt{clients} :
\texttt{medications}, \texttt{allergies}, \texttt{conditions}. Une
nouvelle table \texttt{prescription\_checks} stocke chaque
verification de securite jamais effectuee : le texte de prescription,
le snapshot du contexte patient au moment, le verdict de l'IA, la
decision finale du medecin. C'est la piste d'audit.

La fonction Edge \texttt{prescription-check} prend un
\texttt{client\_id} et une prescription en texte libre, construit un
package de contexte depuis les listes du patient, et appelle Claude
Haiku avec un schema JSON strict couvrant les interactions, les
correspondances allergiques, les precautions, la severite la plus
elevee, et un resume d'une ligne. Le modele est instruit d'etre
pessimiste : en cas d'incertitude, signaler comme precaution plutot
que de laisser passer silencieusement.

\subsection{Decision du medecin et audit}
Le tableau de bord rend le verdict comme une carte coloree --- rouge
pour eleve ou critique, ambre pour moyen, bleu pour bas, vert pour
sur. Trois boutons de decision permettent au medecin d'enregistrer
son choix final : \emph{Annuler la prescription}, \emph{Modifier et
reverifier}, \emph{Outrepasser et continuer quand meme}. Le choix est
horodate sur la ligne \texttt{prescription\_checks}.

Cette piste d'audit est la defense contre la faute professionnelle.
Si un patient poursuit plus tard, la ligne montre que l'IA a signale
le risque, que le medecin l'a vu, et qu'il a soit modifie la
prescription, soit l'a outrepassee avec une justification documentee.

\section{Profils de vehicules pour garages}

\subsection{Le probleme metier}
Les ateliers de mecanique vivent dans un monde different des
cliniques. Leur unite de rendez-vous n'est pas un creneau de
consultation mais une \emph{visite par vehicule} : un changement
d'huile est un travail de quinze minutes, un remplacement de
plaquettes de frein deux heures, un diagnostic sur un bruit non
specifie peut aller de vingt minutes a un apres-midi entier. Chaque
vehicule a sa propre histoire que le mecanicien doit se rappeler :
VIN, kilometrage, dernier changement d'huile, quelles pieces ont ete
utilisees, quelles photos de degats ont ete prises la derniere fois.

\subsection{La solution}
Pour les locataires garages, nous avons cree trois nouvelles tables :
\texttt{vehicles} (VIN, plaque, marque/modele/annee, carburant,
transmission, kilometrage actuel, kilometrage du dernier et prochain
changement d'huile, dates d'expiration assurance et controle
technique), \texttt{vehicle\_services} (ordres de travail
chronologiques avec JSON de pieces, totaux pieces/main d'oeuvre,
kilometrage au service, photos, technicien), et
\texttt{vehicle\_photos} (chronologie de photos separee pour que les
photos WhatsApp casuelles puissent s'attacher sans avoir besoin d'un
enregistrement de service complet).

\subsection{Auto-attache des photos WhatsApp}
La partie la plus interessante est le cote WhatsApp. Lorsqu'un client
envoie une photo avec une legende comme \emph{\og ma Clio fait ce
bruit \fg{}}, le webhook detecte l'image, la telecharge, l'uploade
vers un bucket de stockage \texttt{vehicle-photos}, puis appelle une
fonction Postgres \texttt{find\_vehicle\_for\_message} qui essaie
d'extraire un motif de plaque tunisienne (n'importe quel groupe de 6 a
9 chiffres contigus, avec infixe \texttt{TUN} optionnel) depuis le
texte de la legende et fait correspondre exactement contre
\texttt{vehicles.license\_plate}. En l'absence de correspondance de
plaque, elle essaie un regex VIN de 17 caracteres. Si toujours pas de
correspondance, elle retombe sur le vehicule le plus recemment
enregistre du client.

\section{Surprises d'anniversaire IA}

\subsection{Le probleme metier}
Un client retenu vaut bien plus qu'un client nouvellement acquis,
particulierement dans les entreprises de service ou le bouche-a-oreille
guide l'acquisition. Un petit geste personnalise le jour de
l'anniversaire du client --- un add-on de 30 minutes gratuit, une
remise, un cafe --- est le genre de chose qui convertit un client
satisfait en defenseur vocal. Mais personne ne se souvient des
anniversaires de tout le monde.

\subsection{La solution}
Un cron quotidien \texttt{birthday-runner} a 9h00 scanne chaque
entreprise dont \texttt{birthday\_config.enabled} est vrai. Pour
chaque client dont l'anniversaire correspond a aujourd'hui (toute
annee), le cron frappe un code unique a 6 caracteres depuis un
alphabet de 32 caracteres (sans I, O, 0, 1, tous ambigus), insere une
ligne \texttt{birthday\_vouchers}, et envoie un message WhatsApp dans
la langue du client.

\subsection{Auto-redemption via WhatsApp}
La mecanique virale : quand un client arrive a l'entreprise et envoie
le code nu de 6 caracteres (juste le code, seul), le webhook
court-circuite --- appelle \texttt{find\_birthday\_voucher} pour
valider, appelle \texttt{redeem\_birthday\_voucher} pour marquer
redime, poste une notification proprietaire, et repond
\emph{\og Voucher G7K3MP redime --- Add-on gratuit de 30 minutes.
Profitez ! \fg{}}.

\subsection{Capture de l'anniversaire}
La question implicite : comment Nova connait-il l'anniversaire en
premier lieu ? Trois voies convergent : saisie manuelle via la fenetre
modale Ajouter Client a l'admission ; bouton \og + ajouter \fg{} en
ligne a cote de chaque anniversaire vide dans la table Clients ;
capture conversationnelle (apres une reservation confirmee, si la
fonctionnalite est activee et que le client n'a pas d'anniversaire
enregistre, Nova demande une fois : \emph{\og PS : quelle est votre
date de naissance ? Nous avons une petite surprise pour vous le jour
de votre anniversaire --- ex. 1990-05-23, ou ignorez. \fg{}}). La
troisieme voie est ce qui fait croitre la liste d'anniversaires sans
travail manuel.

\section{Check-in QR et minuteur de consultation}

\subsection{Le probleme metier}
La fonctionnalite d'apprentissage de duree repose sur des horodatages
\texttt{arrived\_at} et \texttt{completed\_at} precis. Les boutons
\emph{Arrive}/\emph{Termine} du tableau de bord fonctionnent, mais
exigent que le secretaire ou le medecin se souvienne de les cliquer,
ce qu'ils oublient souvent.

\subsection{La solution}
Chaque rendez-vous, a l'insertion, recoit automatiquement un
\texttt{checkin\_token} aleatoire de 12 octets encodes en base64 via
un trigger base de donnees. Le token est inclus dans le message de
confirmation WhatsApp comme un lien :
\emph{\og Votre code de check-in (a montrer a l'arrivee) :
https://nova-ai-s8i6.vercel.app/pwa/check.html?a=... \fg{}}

\subsection{Flux de scan medecin}
A la clinique, le medecin (ou proprietaire) ouvre le tableau de bord
et tape l'icone QR violette dans la barre superieure. Une fenetre
modale s'ouvre avec le flux webcam (\texttt{html5-qrcode} charge a la
demande depuis CDN) et un champ de saisie manuelle de token comme
solution de secours.

\subsection{Minuteur en direct}
A un demarrage reussi, le tableau de bord affiche une carte indigo
en haut de page avec le nom du patient en gras 24 pixels, le service
et le telephone en texte attenue, l'horodatage de demarrage, un
minuteur a chiffres tabulaires de 48 pixels comptant en direct dans
le gradient de marque, et un grand bouton rouge \emph{Terminer la
consultation} a droite.

Le panneau de minuteur survit a la navigation entre les pages du
tableau de bord --- l'etat reside dans un objet
\texttt{DASH\_CT\_STATE} de portee module, le handle
\texttt{setInterval} est preserve, et \texttt{bindDashboard} restitue
le panneau a chaque changement de page. Le medecin peut parcourir le
calendrier au milieu d'une visite sans perdre le minuteur.

\section{File d'attente en temps reel (laboratoire)}

\subsection{Le probleme metier}
Un laboratoire medical tunisien typique le matin est chaotique. Les
patients arrivent entre 7h00 et 9h00 pour des prises de sang a jeun,
tous tenant une prescription papier, tous encombrant la petite
reception. Le laboratoire ne peut pas les servir simultanement, ils
font donc la queue. La salle d'attente devient saturee, les patients
perdent la trace de leur position, le secretaire est constamment
interroge sur \emph{\og combien de temps ? \fg{}}.

\subsection{La solution}
Chaque locataire laboratoire recoit un \emph{QR de reception}
imprimable qui encode une URL pointant vers
\texttt{/pwa/queue.html?b=<biz-id>}. Le proprietaire l'imprime sur du
A4 et le colle sur le bureau.

Un patient arrivant au laboratoire scanne le QR avec la camera de son
telephone, atterrit sur une page a carte unique, tape son nom et
numero de telephone, et tape \emph{Obtenir ma position}. La RPC
\texttt{queue\_join} insere une ligne \texttt{queue\_entries},
retourne le nouveau numero de position et un temps d'attente estime
calcule depuis la moyenne de temps de service des 20 dernieres
entrees.

\subsection{Mises a jour en direct}
La carte du patient se transforme en une vue de statut : un numero
de position de 96 pixels dans le gradient de marque, un ETA de 32
pixels (\emph{\og environ 22 min \fg{}}), et une sous-ligne expliquant
la moyenne. Une souscription Supabase Realtime sur la table
\texttt{queue\_entries} filtree par \texttt{business\_id} pousse les
mises a jour chaque fois qu'une entree dans la file change. La
position du patient se recalcule instantanement ; il peut quitter le
batiment, prendre un cafe, et revenir quand sa position passe a
\og C'est votre tour \fg{}.

\section{Avatars vocaux}
Le clonage vocal a ete implemente mais conserve derriere un drapeau
de fonctionnalite (\texttt{businesses.voice\_feature\_enabled}) en
attendant un abonnement ElevenLabs paye. Le pipeline technique est en
place : le MediaRecorder dans le navigateur capture un echantillon de
30 secondes, la fonction Edge \texttt{voice-avatar} l'uploade vers
ElevenLabs, stocke l'identifiant de voix retourne, et le webhook
WhatsApp attache une note audio text-to-speech a cote de chaque
reponse texte quand la fonctionnalite est activee.

\section{Revue du sprint et tests}
Les tests de bout en bout du Sprint 3 ont dure une semaine complete
avec trois entreprises reelles (une clinique, un laboratoire, un
garage). Faits saillants : la clinique a execute une campagne
demand-fill pour un vendredi apres-midi vide. Sur douze patients
pinges, quatre ont repondu OUI, deux ont reserve. Taux de remplissage :
16,7\%. Le laboratoire a utilise la file d'attente en temps reel
pendant trois matinees consecutives ; le secretaire a rapporte que la
baisse de saturation de la salle d'attente etait notable. Le garage a
teste l'auto-attache de photo avec cinq photos, quatre ont fait
correspondance correctement via la plaque, une via le repli. Le
verificateur de securite des medicaments a signale deux interactions
reelles dans vingt prescriptions de test, toutes deux correctement.

\section{Conclusion}
Le Sprint 3 a clos la boucle de developpement. Avec les huit
fonctionnalites d'intelligence en place, Nova~AI n'est plus juste un
outil de prise de rendez-vous automatise --- c'est un partenaire
operationnel quotidien qui sauve proactivement le proprietaire des
modes de defaillance courants (creneaux vides, consultations qui
debordent, prescriptions dangereuses, anniversaires oublies, historique
vehicule perdu). Le chapitre final tire les conclusions et trace les
directions d'evolution future que ce travail a ouvertes.
